\section{Background}\label{ch1:sec:background}

%introduction on 3D modeling and reconstruction
Geometric models are the kernels of many scientific and engineering
subjects. They are widely used in various applications including
product design, digital entertainment, biology, archeology and so
on~\cite{BPKALBR07}. Correspondingly, there is a strong need for
development of modeling and reconstruction techniques to create,
edit, analyze and compute 3D models. Although there
have been a lot of techniques and devices such as image-based modeling~\cite{DTM96,LTZYWK06} and 3D scanner techniques~\cite{LPCRK00,ZSCS04}
developed to build virtual 3D models, most 3D shapes still need to
be constructed manually via various modeling tools by the user,
either from scratch or by editing existing models.

\begin{figure} [htbp]
    \centering
  \includegraphics[scale=0.35]{figs/f1.geomodels.eps}
  \caption{Geometric models used in various applications. These models are courtesy of the INRIA Gamma team research database~\cite{INRIA}.}
\end{figure}

Traditional tools for 3D modeling and reconstruction such as 3ds Max~\cite{3DSMAX}, Maya~\cite{MAYA} and so on are mainly
designed for expert users to precisely construct sophisticated
models~\cite{CIW08}. They are difficult to use. One must spend a long time
learning the tools and is expected to work on a task for a period of
time to obtain a final result. This seems too complicated for novice
users to master. These tools are also difficult for professional
users in the initial design phase. Designers still need to work with
pen and paper first and then move on to the computational tools when
the design is fixed.

\begin{figure} [htbp]
    \centering
  \includegraphics[scale=0.2]{figs/f1.3dsMax.eps}
  \caption{3D modeling using the 3ds Max software~\cite{3DSMAX}.}
\end{figure}

Advances in information technology have  enabled many systems to
extend human's abilities in design and creative work. Modern
cognition study indicates that human has two cognitive modes:
experimental cognition and reflective cognition~\cite{DN98}. The
former corresponds to skilled behavior and the latter accounts for
intellectual work. Tano et al.~\cite{TKNK03} show that while most of
current computer graphics (CG) and virtual reality (VR) tools are
good for experimental cognition and poor for reflective cognition,
sketch-based tools can promote reflective cognitive activities. On
the other hand, with the rapid development of pen-based computer
hardware, digital sketching interface emerges as a powerful way for
users to interact with the computer and do various 3D modeling and
reconstruction tasks.

Sketch-based modeling refers to the process  of producing 3D models
from user sketched curves or strokes. Instead of requiring the users
to work with buttons, menus, icons and pointers, it allows the
direct expression of their thoughts in the form of freehand
sketching. The goal of the sketching interfaces is to solve the
problem of accessibility by reducing the amount of explicit control
of the users to simplify the user interface dramatically. The
sketch-based technique makes 3D modeling and reconstruction tasks
accessible to novice users. Meanwhile, it can also be used by
experts in their initial design stage. With the improvement of the
computational capabilities of computers and the emerging of various
powerful input devices, the sketching techniques are developing
rapidly in recent years. There have been some commercial modeling
software and packages supporting freehand sketches, examples of
which are SketchUp~\cite{SketchUp}, Archipelis
Designer~\cite{ArcDes}, and Sunny 3D~\cite{Sunny3D}.

\begin{figure} [htbp]
    \centering
  \includegraphics[scale=0.6]{figs/f1.sketch.eps}
  \caption{Sketch-based 3D modeling.}
\end{figure}

The design of interactive and  powerful sketching interfaces
requires the exploration of various modeling algorithms and the
combination of the knowledge from other fields, such as human
psychology, virtual reality, computer vision and so on. At the same
time, the sketching technique can also help to make the fundamental
modeling and reconstruction algorithms more effective and
intelligent. The development of sketch-based modeling systems and
the utilization of the sketching technique in improving the modeling
and reconstruction algorithms compose the main content of this
research.


\section{Problem statement and objectives}

Though a lot of efforts have been made in sketch-based modeling and
reconstruction, there are still many important open problems.
Sketch-based modeling is often related to human shape perception and
human shape perception of 2D sketches is a complicated issue
involving various subjects. While sketch-based techniques offer an
improvement over traditional methods in terms of accessibility, they
are not yet complete replacements. There remains much work to be
done for sketch-based modeling  to produce a wide range of objects
with high complexity. The challenges for sketch-based modeling
include the development of a natural and friendly user interface,
the exploration on perceptually meaningful description of a 3D
shape, the proper interpretation of the user sketches, and the
enhancement of the sketch precisions and model qualities. Meanwhile,
it is non-trivial to incorporate sketching techniques into
fundamental modeling and reconstruction algorithms in order to
produce a faithful and desirable results to the user.


This thesis researches the use  of sketching techniques in 3D
modeling and reconstruction to enhance various geometric modeling
algorithms. It aims to gain deep understanding of diverse techniques
from geometric modeling, computer graphics, human computer
interaction required to advance the state of the art in interactive
3D modeling, to investigate how sketching interfaces can be
efficiently used for enhancing 3D modeling and reconstruction, and
to develop sketch-based techniques for creating, editing, and
manipulating digital geometric models. In particular, the research
objectives are:
\begin{itemize}
\item to explore and develop novel sketching interfaces and
algorithms for sketch-based modeling, by combining the knowledge of
geometric modeling, human computer interaction, compute vision, human psychology and so on;

\item to investigate how to enhance the fundamental 3D modeling and reconstruction
algorithms by using sketching techniques, thus producing faithful
and satisfactory modeling results to the user and making the
algorithms suitable for various interactive geometric
 processing applications.
\end{itemize}


\section{Contributions and thesis organization}


%approaches:
To accomplish these objectives, the following approaches have been
proposed:

%system: reference plane and rules
First, we developed a sketch-based modeling system with auxiliary
planes as  references for 3D freeform shape design. The user can
first create an initial 3D shape with arbitrary topology by
sketching the cross section curves of the model, and then use
sketches to perform operations such as deformation, extrusion and so on
to further edit the
model. To regularize and interpret the user's sketches properly, we
introduce some rules for the input strokes, which are based on both
the semantic meaning of the sketched strokes and human psychology.
Unlike other sketching systems, all the creation and editing
operations in the proposed system are performed with reference to
some auxiliary planes that are automatically constructed based on
the user sketches or default settings. The use of reference planes
provides a heuristic solution to the ambiguity problem of the 2D
interface for modeling in 3D space and makes the user input
meaningful and intuitive.

%progressive approach and algorithm, to enhance the creation function
Second, it is observed that most previous sketch-based modeling
systems mainly focus on providing editing tools to gradually sculpt
a 3D model from a rough shape rather than creating it from scratch
through iterative sketches. In order to enhance the creation
function in sketching-based modeling, we proposed a progressive
modeling approach, which allows the user to create a 3D model
through iterative sketches, taking both existing sketched curves and
the up-to-date reconstructed 3D shape as references to get fully
aware of the shape of the sketched model. A novel surface
reconstruction algorithm which ensures gradual shape updating during
iterative sketching and produces models with a singly connected
component is proposed to effectively support this creation
operation.
%detailed introduction on the algorithm?

%deformation
Third, we studied one important editing function in sketch-based
modeling: surface deformation. We presented a new surface-based mesh
deformation method that performs the computation via an edge-based
graph, which allows natural increase of the sampling rate for more
accurate shape computation and generates better deformation results.
The user is given the flexibility of adjusting the deformation
effect between local shape preservation and global smoothness.
Moreover, to simulate the deformation behaviors of regions with
different materials, we introduce a stiffness property into the
deformation model and present an easy and intuitive way for the user
to set the material property as well. Thus, our algorithm can produce
satisfactory results and make the deformation more flexible.

%surface reconstruction
Fourth, we utilized sketching  in solving the problem of surface
reconstruction from cross sections with arbitrary orientations and
presented a new scheme which constructs a smooth mesh surface from
cross section curves. The scheme is composed of two parts. The first
part is a surface reconstruction algorithm that produces a manifold
surface interpolating the input cross sections and having regular
local shapes. It is the  generalization of the algorithm developed
in our progressive modeling which is used to handle cross sections with arbitrary
orientations. The second part is sketching assisted topology
editing, which allows the user to make further adjustment on the
local topology of the reconstructed model through simple sketches if
the initial shape is not satisfactory. The computation of this
editing operation makes use of the intermediate information obtained
during the reconstruction process and avoids complicated analysis on
the shape of the global surface. Our surface reconstruction scheme
has the unique ability to build a manifold surface with only one
connected component, and avoids producing spurious or weired local
shapes which are often seen in previous methods. The reconstruction
process is further accelerated by computing the sub-surfaces in
parallel. This surface reconstruction scheme can be used in various
applications including sketch-based modeling.


The remainder of this thesis is organized as follows:
\begin{itemize}
\item Chapter~\ref{ch:review} reviews the related works, which include
 sketch-based modeling systems, mesh deformation algorithms and
  surface reconstruction from cross section curves.

\item Chapter~\ref{ch:planeSBIM} presents our sketch-based modeling
system which takes planes as references for freeform shape design.
The interfaces for the progressive creation and editing operations,
the rules for interpreting and regularizing the user sketches, as well as
 the algorithm on the calculation of the reference planes in this system are
 described in detail.
%The system also serves a vehicle for carrying out our research in next steps.

\item Chapter~\ref{ch:orthsurf} introduces the algorithm on progressive
surface reconstruction from orthogonal cross sections, which effectively
 supports the creation function in our sketching interface.

\item Chapter~\ref{ch:flexdeformation} presents the algorithm on
mesh deformation , which is one of the most important editing functions in our system.

\item Chapter~\ref{ch:editsurf} introduces a new framework for
surface reconstruction from cross sections with arbitrary
orientations, which includes generalizing the algorithm in
Chapter~\ref{ch:orthsurf} to handle cross
 sections with arbitrary orientations and using sketching  to edit
 topology of the resulting surface.

\item Chapter~\ref{ch:conclusion} summarizes the thesis and points out some
directions for future work.
\end{itemize}
